\section{Límites de la evidencia y términos clave}

Esta clase abarca arquitectura, sistemas y resultados de producto, por lo que es fácil caer en dos malentendidos: primero, sustituir una evaluación integral por una sola muestra; segundo, sustituir la latencia real por los FLOPs. Las preguntas de la sesión de Q\&A sobre GAN, métricas automáticas, costo del video, AdaLN y MMDiT aportan los límites de ingeniería necesarios.

\subsection{Cuatro correcciones que deja la sesión de Q\&A}

Primero, las GAN no han desaparecido por completo: siguen siendo valiosas cuando la generación one-shot y la latencia ultrabaja importan más que la máxima calidad. Segundo, las métricas automáticas de imagen y video son débiles; deben combinarse con la preferencia humana y con métricas de la tarea. Tercero, un video no es un conjunto de cuadros independientes. Cuarto, la intuición detrás de MMDiT no es una demostración de optimalidad teórica.

\teachervoice{00:57:31--01:01:38, la sesión de preguntas en línea trata los modelos autoregressive, las GAN, las métricas automáticas y los retos del video. El docente no responde que «la difusión gana en todos los casos», sino que distingue los escenarios según latency, quality y temporal consistency.}

\teachervoice{01:01:45--01:03:00, el docente vuelve a explicar AdaLN: LayerNorm se encarga de estabilizar las características y la parte adaptive usa la condición externa para modificar scale/shift, de modo que el mismo block ejecuta cálculos distintos según el timestep o la condición.}

\begin{warningbox}{Qualitative cherry-picking}
Un modelo generativo permite elegir, entre una gran cantidad de muestras, unas pocas imágenes óptimas. Un reporte confiable debe fijar el prompt set, publicar las seeds, mostrar las muestras fallidas y medir al mismo tiempo prompt fidelity, diversity, human preference, latency y memory.
\end{warningbox}

\subsection{Concentrado de términos clave}

A continuación se reorganizan los términos de alta densidad del curso según «el problema que resuelven», para no confundir el nombre de un modelo con un mecanismo.

\begin{center}
\small
\begin{tabularx}{\textwidth}{p{0.18\textwidth} X X}
\toprule
\textbf{Término} & \textbf{Problema que resuelve} & \textbf{Mecanismo central y relación con el curso} \\
\midrule
Scheduler & Cómo pasar de un estado de ruido al siguiente & Define el timestep, el programa de ruido y la actualización numérica; no equivale al denoiser. \\
Denoiser & Qué dirección debe predecir en cada paso & U-Net o DiT; produce una parametrización de noise, velocity o data. \\
Latent diffusion & Cómo reducir el costo del espacio de píxeles & Primero comprime con un VAE y luego itera sobre un latent de baja resolución. \\
U-Net & Cómo aprovechar el sesgo multiescala de las imágenes & Los down/up blocks y las skip connections conservan la información local y global. \\
DiT & Cómo eliminar ruido con un Transformer estándar & Aplica patchify al latent y modula los encoder blocks con timestep/condition. \\
AdaLN-Zero & Cómo inyectar la condición a bajo costo y estabilizar redes profundas & La condición genera scale/shift/gates; la inicialización en cero hace que el block sea inicialmente casi la identity. \\
Cross-attention & Cómo leen el texto los token de imagen & La imagen funciona como query y el texto como key/value; normalmente solo se actualiza el flujo de imagen. \\
Linear attention & Cómo evitar el $N^2$ de la image self-attention & Reordena la secuencia de multiplicaciones de la attention kernelizada, a costa de parte de la interacción exacta. \\
MMDiT & Cómo conservan las dos modalidades su espacio propio y evolucionan juntas & AdaLN/QKV separados; tras concatenar, se aplica joint attention. \\
QK-Norm & Cómo estabilizar los attention logits en modelos grandes & Normaliza query/key para controlar la escala del producto punto. \\
GQA & Cómo reducir el costo de key/value & Varias query heads comparten menos K/V heads. \\
RoPE & Cómo codificar la posición relativa espacial y temporal & Rota query/key; puede extenderse a los tres ejes del video. \\
Structural control & Cómo seguir edge/pose/depth & Codifica condiciones alineadas espacialmente y las inyecta en el base model. \\
Parameter sharing & Cómo reducir los pesos y parte del cálculo & Comparte AdaLN, QKVO y MLP entre capas o entre modalidades, con un compromiso entre calidad y eficiencia. \\
\bottomrule
\end{tabularx}
\end{center}

\subsection{Resumen de la sección}

Las arquitecturas generativas no pueden compararse sin tomar en cuenta el objetivo y la workload. GAN, U-Net, DiT, linear attention y MMDiT tienen cada una su rango de aplicación; las métricas automáticas, el número de parámetros y los FLOPs son solo indicadores parciales. Una conclusión confiable requiere control de variables, mediciones de extremo a extremo y muestras fallidas.
